\documentclass[10pt,twocolumn,letterpaper]{article}

% --- Page geometry ----------------------------------------------------------
\usepackage[letterpaper,margin=0.75in,columnsep=0.28in,top=0.85in,bottom=0.95in]{geometry}

% --- Encoding & fonts -------------------------------------------------------
\usepackage[utf8]{inputenc}
\usepackage[T1]{fontenc}
\usepackage{lmodern}
\usepackage{microtype}

% --- Common packages --------------------------------------------------------
\usepackage{graphicx}
\usepackage{booktabs}
\usepackage{tabularx}
\usepackage{array}
\usepackage[dvipsnames]{xcolor}
\usepackage[most]{tcolorbox}
\usepackage{caption}
\usepackage{amsmath,amssymb}
\usepackage{enumitem}
\usepackage{url}
\usepackage{seqsplit}
\usepackage{fancyvrb}
\usepackage{fvextra}
\usepackage{titlesec}
\usepackage{abstract}
\usepackage[colorlinks=true,linkcolor=black,citecolor=NavyBlue,urlcolor=NavyBlue]{hyperref}

% --- Tighter section spacing ------------------------------------------------
\titlespacing*{\section}{0pt}{1.4ex plus .5ex minus .2ex}{0.6ex plus .2ex}
\titlespacing*{\subsection}{0pt}{1.2ex plus .4ex minus .2ex}{0.5ex plus .2ex}
\titleformat*{\section}{\large\bfseries}
\titleformat*{\subsection}{\normalsize\bfseries}
\titleformat*{\subsubsection}{\normalsize\itshape}

% --- Tighter list spacing ---------------------------------------------------
\setlist{itemsep=2pt,topsep=3pt,parsep=0pt}

% --- Caption styling --------------------------------------------------------
\captionsetup{font=small,labelfont=bf,skip=4pt}

% --- Title formatting tweaks (centered, no rule, tighter spacing) -----------
\makeatletter
\renewcommand{\@maketitle}{%
  \newpage
  \null
  \begin{center}%
  {\@title \par}%
  \vspace{0.6em}%
  {\@author \par}%
  \vspace{0.4em}%
  {\@date \par}%
  \end{center}%
  \par
  \vspace{0.6em}%
  \rule{\textwidth}{0.4pt}%
  \vspace{0.2em}%
}
\makeatother

% --- Abstract title smaller and tighter -------------------------------------
\renewcommand{\abstractnamefont}{\normalfont\small\bfseries}
\renewcommand{\abstracttextfont}{\normalfont\small}
\setlength{\absleftindent}{0pt}
\setlength{\absrightindent}{0pt}

% --- Big-picture callout box ------------------------------------------------
\newtcolorbox{takeaway}[1][]{
  enhanced, breakable, sharp corners=south,
  colback=blue!4!white, colframe=blue!50!black,
  boxrule=0.6pt, left=6pt, right=6pt, top=4pt, bottom=4pt,
  fonttitle=\bfseries, title={Takeaway}, #1
}

% --- Inline code and small numerals -----------------------------------------
\newcommand{\code}[1]{\texttt{\small #1}}
\newcommand{\sha}[1]{\texttt{\footnotesize\seqsplit{#1}}}

% --- Allow LaTeX to stretch to avoid overfull boxes -------------------------
\setlength{\emergencystretch}{2.5em}
\tolerance=1500

% --- Title block ------------------------------------------------------------
\title{\vspace{-0.4em}\bfseries\Large
  Dynamic / ME-JEPA\\[3pt]
  \large An Audited, Domain-Portable World-Model Runtime\\
  on a Single RTX 5090, with a New Class of Training Data
  \vspace{0.2em}}

\author{%
  \normalsize Chris Royse\\[1pt]
  \footnotesize Independent Researcher\quad\textbullet\quad
  \texttt{chrisroyseai@gmail.com}%
}

\date{\footnotesize May 2026 \quad\textbullet\quad
  Release candidate
  \texttt{mejepa\_5090\_artifact\_v1.0.0\_rc5\_sif\_verified}}

% ============================================================================
\begin{document}

% --- Banner: title + abstract span the full page width ---------------------
\twocolumn[
  \maketitle
  \begin{onecolabstract}
  \noindent
  Joint-embedding predictive architectures (JEPA) learn to predict
  latent representations of the future instead of reconstructing
  pixels or tokens. Most JEPA work focuses on representational
  quality on natural images or video. This paper takes a different
  angle: \emph{what does it take to make every step of a JEPA
  runtime --- data ingest, panel materialization, action-conditioned
  training, planning, surprise recording, and artifact provenance ---
  physically auditable on a single commodity GPU?}\par\vspace{3pt}\noindent
  %
  Dynamic / ME-JEPA is one Rust binary that does all of it. A new
  domain ships only as a TOML manifest. Every operation writes typed
  rows into 25 RocksDB column families plus an append-only audit log,
  and a structurally separate verifier reads everything back. The
  same code path runs three domains: \texttt{counter\_world}
  (480 trainable parameters), \texttt{gridworld\_5x5} (16{,}960), and
  a new public-taxonomy career benchmark
  \texttt{career\_taxonomy.v1} (15{,}344) that walks the related-
  occupation graph of O*NET 30.2 starting from Software Developers
  \texttt{15-1252.00}. Five seeds of \texttt{paper\_small\_1k}
  (1{,}000 events each) on an RTX 5090 give career validation cosine
  0.905~$\pm$~0.044 (BCa 95\% CI 0.861--0.931), test cosine
  0.868~$\pm$~0.068, and validation action-contrast accuracy
  0.960~$\pm$~0.043. The negative-action ablation A3 collapses to
  validation cosine 0.790, exactly as the design predicts and the
  verifier confirms.\par\vspace{3pt}\noindent
  %
  The second contribution is conceptual. Each immutable bundle
  produced by the runtime is itself a structured training-data
  artifact. Most of its bytes are deterministic functions of a
  public source ZIP; the small model-output slice is paired with the
  observed next-state and with random-init / shuffled-target
  baselines in the same record; counterfactual ablations and failed
  runs are persisted as first-class records. This paper argues that
  bundles form a new class of training data --- \emph{audited
  execution traces with cryptographic provenance} --- distinct from
  web text, pure model-generated synthetic data, and rejection-
  sampled chain-of-thought, with a substantially lower
  recursion-replacement collapse hazard. The reproducibility
  artifact is a local Docker + Singularity (SIF) release candidate
  with deterministic Docker tar SHA-256 \sha{099f053\dots7d4cd80e},
  SIF SHA-256 \sha{30b531c\dots f2825}, archive SHA-256
  \sha{8cccf49\dots43ecd}.

  \vspace{4pt}
  \noindent\textbf{Keywords.} JEPA, world models, auditable ML,
  reproducibility, training data, model collapse, process
  supervision, RTX 5090.
  \end{onecolabstract}
  \vspace{6pt}
  \rule{\textwidth}{0.4pt}
  \vspace{8pt}
]

% ============================================================================
\section{Introduction}
\label{sec:intro}

Imagine that someone hands you a paper claiming a new latent
world-model trains on a single GPU and beats baselines on a career-
domain task. You want to reproduce it. You go through the motions:
download the model weights, find the dataset (somewhere), spin up
a few-machine cluster, hope your CUDA versions match, and then
realize that the path from raw event $\to$ training row $\to$
predictor weights $\to$ planning trace $\to$ verified output is not
written down anywhere as a single deterministic pipeline. You can
\emph{rerun} the model. You cannot \emph{prove} the entire pipeline
behaves the way the paper says it does.

That gap is the motivation for this work.

This paper is \textbf{not} a representational-quality contribution
to the JEPA family. The premise of latent prediction is taken as
established \cite{assran2023ijepa,bardes2024vjepa,lecun2022path}.
The question is one level up:
\begin{quote}
\itshape
What would it take to make a complete latent world-model runtime ---
data ingest, multi-instrument latent panel materialization,
action-conditioned training, planning, surprise recording, and
artifact provenance --- physically auditable end-to-end on a single
commodity GPU?
\end{quote}

\textbf{Dynamic / ME-JEPA} is one Rust binary that answers it.
Twenty-nine subcommands compose into one immutable bundle per run.
A new domain ships as a TOML manifest. Twenty-five typed RocksDB
column families plus an append-only audit log capture every
persistent state. A structurally separate verifier code path reads
everything back. Three domains are instantiated by domain pack
alone, including a new public-taxonomy career benchmark grounded in
O*NET 30.2 \cite{onet2026}. Five seeds with bias-corrected and
accelerated bootstrap confidence intervals
\cite{efron2020automatic} and paired sign-flip permutation tests
give the headline numbers in \S\ref{sec:results}. A
deterministically built Docker + Singularity (SIF) artifact ships
with full SHA-256 chains in \S\ref{sec:repro}.

The second contribution is conceptual. After running the system,
each immutable bundle is itself a small structured corpus. The
deterministic walk over O*NET, the materialized panels, the audit
log, the typed edge-case rows, the ablation rows, and the failed-
run records make a bundle a candidate \emph{training-data artifact}
in its own right. \S\ref{sec:training-data} argues that bundles
form a new class of training data and that this class carries lower
recursion-replacement collapse hazard \cite{shumailov2024nature}
than typical synthetic data because most of the bytes are
deterministic functions of a public source rather than model
outputs.

\subsection{Why a Domain-Portable Runtime}

The dominant pattern in JEPA papers is a representational claim
backed by a custom training stack. Useful for the claim. Hard to
extend. Each new domain rebuilds the audit and provenance
infrastructure from scratch, if it builds them at all. The
domain-pack abstraction in \S\ref{sec:method} makes that
infrastructure once and reuses it everywhere a TOML manifest can
be written.

\subsection{Why a Public-Taxonomy Career Benchmark}

The career-domain has the property that data pipelines benefit from
provenance recorded inline, even when the model is small. O*NET
30.2 is the canonical public occupation taxonomy in the United
States, distributed under CC BY 4.0 with explicit attribution
requirements. Constructing a benchmark that walks O*NET's related-
occupation graph deterministically gives a substrate where every
generated transition can be traced back to a public source row by
SHA-256 hash. Section~\ref{sec:bench} specifies the benchmark.

\subsection{Why Single-GPU Feasibility}

Restricting compute to a single host does two useful things at
once. It removes the "hidden cluster" reproduction barrier that
makes most public ML results expensive to verify, and it forces
training-budget honesty: a five-seed paper-small ladder plus
container build plus container verify plus manual edge-case audits
fits in roughly two hours of wall time on the RTX 5090 used here.

\subsection{Contributions}

In order:
\begin{enumerate}
  \item \textbf{A domain-pack-driven JEPA runtime.} One Rust binary
        runs three domains by TOML manifest alone, with no domain-
        specific code branches (\S\ref{sec:method}, App.~A).
  \item \textbf{The \texttt{career\_taxonomy.v1} public benchmark.}
        Generated deterministically from O*NET 30.2 starting from
        \texttt{15-1252.00}, with per-edge walk hashes and an
        archive preflight (\S\ref{sec:bench}).
  \item \textbf{A strict full-state verification protocol.} Every
        claim in \S\ref{sec:results} has a corresponding RocksDB
        row, file SHA-256, and decoded readback through a separate
        code path (\S\ref{sec:fsv}).
  \item \textbf{A new training-data class: audited execution traces
        with cryptographic provenance.} Bundles are structured
        corpora of deterministic-from-source records, tiny-model
        outputs paired with observed targets, and independent
        verifier readbacks. Counterfactual ablations and failed
        runs are first-class records (\S\ref{sec:training-data}).
  \item \textbf{Multi-seed paper-small results with BCa CI and
        paired permutation tests.} Five-seed career validation
        cosine $0.905 \pm 0.044$ with BCa 95\% CI \([0.861, 0.931]\);
        every cited number reads from a CSV in the verified release
        bundle (\S\ref{sec:results}).
  \item \textbf{A reproducibility artifact.} Local Docker + SIF
        release candidate with full hash chain
        (\S\ref{sec:repro}).
\end{enumerate}

\subsection{Scope}

The trained models are deliberately small: $480$, $16{,}960$, and
$15{,}344$ trainable parameters. The point is not "small model
beats large model"; it is that every stage of a portable JEPA
system can be physically verified on commodity hardware, and that
this verification produces a structured corpus as a side effect.

% ============================================================================
\section{Related Work}
\label{sec:related}

This work sits at the intersection of four lines of literature.

\subsection{Joint-Embedding Predictive Architectures}

I-JEPA \cite{assran2023ijepa} predicts the representations of
target image blocks from a context block, in a self-distilled
latent space. V-JEPA \cite{bardes2024vjepa} extends the idea to
video by feature prediction, with downstream evaluations showing
that latent prediction recovers strong representations without
reconstruction, negatives, or text supervision. The official Meta /
FAIR V-JEPA repository is the reference engineering implementation
\cite{metafair2024vjepa}. LeCun's "A Path Towards Autonomous
Machine Intelligence" frames the family as a step toward
predictive systems that reason in abstract representation space
\cite{lecun2022path}. Recent JEPA work continues to refine the
training-stability and modality-coverage axes: LeWorldModel
\cite{maes2026leworldmodel} demonstrates stable end-to-end
JEPA training from raw pixels with a single regularizer, and
VL-JEPA \cite{chen2026vljepa} extends the predictive-latent
formulation to vision-language with parameter-efficiency gains
over token-space objectives. The Dynamic / ME-JEPA contribution
is orthogonal to both: it does not propose a new representation
objective, and instead asks what infrastructure is required to
make a complete JEPA runtime --- including data, training,
planning, and provenance --- physically auditable end-to-end.

The Dynamic / ME-JEPA architectural change is in \emph{how target
representations are sourced}: instead of an exponential-moving-
average target encoder, ME-JEPA uses a panel of frozen, externally-
defined per-domain instruments declared by the domain pack
(\S\ref{sec:method}). The narrower claim is that external frozen
targets eliminate target-encoder collapse but \emph{do not}
eliminate predictor degeneracy --- so VICReg-style diagnostics
\cite{bardes2022vicreg} and explicit ablations are still required.

Self-supervised pre-cursors that inform the regularization choices:
BYOL \cite{grill2020byol} (asymmetric online/target with
stop-gradient as collapse mitigation) and Contrastive Predictive
Coding \cite{vandenoord2019cpc} (latent forecast objective).

\subsection{World Models and Latent Planning}

The planning, surprise, and guard plane is informed by World Models
\cite{ha2018worldmodels}, the Dreamer line including DreamerV3
\cite{hafner2023dreamerv3}, and MuZero \cite{schrittwieser2019muzero}.
The imagined-rollout idea conceptually informs the candidate-action
planner (\S\ref{sec:method}), but rollouts are constrained to one
step in the present runtime; multi-step ($H{>}1$) rollouts are an
explicit deferred axis (\S\ref{sec:limits}).

\subsection{Generalist and Multi-Modal Architectures}

Gato \cite{reed2022gato} demonstrated multi-modal multi-task multi-
embodiment tokenized policies; Perceiver IO
\cite{jaegle2022perceiverio} showed that one attention-based
architecture handles arbitrary structured inputs and outputs.
Dynamic / ME-JEPA borrows neither as engineering, but inherits the
constraint they imposed: a general-purpose latent runtime must not
hard-code modality. The domain-pack abstraction is the answer.

\subsection{Career and Occupational Taxonomies}

The U.S. Department of Labor and Employment Training
Administration's O*NET database \cite{onet2026} is the canonical
public occupation taxonomy used here, currently at release 30.2,
distributed under CC BY 4.0. The European Skills, Competences,
Qualifications and Occupations classification (ESCO)
\cite{esco2025} is the EU counterpart and a planned cross-walk
replication source. Recent career-trajectory dataset papers ---
KARRIEREWEGE \cite{senger2024karrierewege} and JobHop
\cite{johary2025jobhop} --- are relevant as future extensions; the
only ingested source in the present runtime is O*NET 30.2.

\subsection{Auditable ML Systems}

The runtime's evidence-bundle design is informed by MLflow Tracking
\cite{mlflow2024} (run, experiment, artifact, metric lineage). The
reproducibility ladder --- pinned base image digest, deterministic
build, deterministic tar export, container DOI --- follows
Docker reproducible-build guidance \cite{docker2025reproducible} and
Apptainer SIF-from-docker-archive guidance
\cite{apptainer2025buildsif}. PyTorch's reproducibility notes
\cite{pytorch2026reproducibility} are inherited as a caveat: exact
GPU result reproducibility across releases or platforms is not
guaranteed; cross-run comparisons therefore use structural plus
tolerance checks, not byte-equal model-file diffs.

\subsection{Synthetic Data, Model Collapse, Process Supervision}

Recent work establishes a non-trivial taxonomy of training-data
classes and their effect on model quality under recursive
training. Shumailov et al.\ introduced \emph{model collapse}:
training generative models on recursively generated data leads to
irreversible loss of distribution tails and eventual convergence to
a degenerate distribution \cite{shumailov2023curse,
shumailov2024nature}. Gerstgrasser et al.\
\cite{gerstgrasser2024accumulate} subsequently showed that
\emph{accumulating} real and synthetic data avoids the collapse
pattern that \emph{replacement} triggers --- distinguishing the
dangerous recursion-with-replacement case from the safer
accumulation case.

In parallel, the highest-quality synthetic reasoning corpora have
moved away from pure LLM generation toward verifier-grounded data:
process supervision and the PRM800K dataset of step-level labels
\cite{lightman2023verify}; code-execution-grounded supervision
\cite{jung2025codeexec}; SYNTHETIC-1
\cite{primeintellect2025synthetic1} (1.4M verified reasoning
traces); and Enigmata \cite{seed2025enigmata}.
Section~\ref{sec:training-data} extends this taxonomy with a class
of training data the Dynamic / ME-JEPA runtime produces by
construction.

% ============================================================================
\section{The Dynamic / ME-JEPA Runtime}
\label{sec:method}

The runtime is a single Rust binary, \code{context-graph}, with
the DynamicJEPA subcommand surface listed in
\S\ref{sec:method:cli}. The core abstraction is a typed transition
over a multi-instrument latent panel:
\begin{equation*}
  (s_t,\, a_t,\, c_t,\, g_t)\ \longrightarrow\ \hat{s}_{t+1}
\end{equation*}
where $s_t$ is the panel-encoded state, $a_t$ is a typed action,
$c_t$ is contextual scaffolding, and $g_t$ is an optional goal
panel. The target $s_{t+1}$ is the panel produced by frozen,
externally-defined instruments declared by the domain pack; the
trainable predictor $P_\psi$ produces $\hat{s}_{t+1}$.

\subsection{Domain Packs}

A domain pack is a strict TOML manifest registering a domain into
the DynamicJEPA RocksDB column families. Required fields:
\code{[domain]}, \code{[time]}, repeated \code{[[state.fields]]},
\code{[[action.fields]]}, \code{[[outcome.fields]]}, and
\code{[[entity.fields]]}. Field kinds are
\code{categorical\{variants\}}, \code{i64\{min,max\}},
\code{f64\{min,max\}}, \code{vector\{dim,min,max\}}, \code{string},
and \code{bool}. Validation is strict: out-of-range values, missing
required fields, undeclared categorical variants, and vector
dimension mismatches each produce typed errors with field paths.
There is no fallback or coercion.

Three domain packs ship with the artifact:
\code{counter\_world.v1} (480 trainable parameters; sanity domain),
\code{gridworld\_5x5.v1} (16{,}960 parameters; compositional
next-state prediction), and \code{career\_taxonomy.v1} (15{,}344
parameters at the headline panel; defined in \S\ref{sec:bench}).
The full \code{career\_taxonomy.v1} TOML is reproduced in
Appendix~\ref{app:domainpack}.

\subsection{Storage: 25 Typed Column Families plus Audit Log}

DynamicJEPA stores data in a dedicated RocksDB instance with 25
column families, all prefixed \code{dj\_*}: domain packs,
instrument and adapter registries, raw events, normalized states,
actions, outcomes, transitions, adapter runs, instrument readings,
latent panels, bindings, trajectories, dataset shards, training
runs, model artifacts, predictions, skill policies, plan traces,
guard decisions, surprise events, verification runs, and the
append-only audit log.

Every mutating subcommand writes one or more CFs and \emph{always}
writes a corresponding \code{dj\_audit\_log} row. Reading a CF back
is a separate code path (\code{inspect-counts}, \code{inspect-cf},
\code{get-*}) that opens the database read-only and re-decodes the
bytes. This keeps the writer from being the verifier. Read-only
opens were hardened during smoke-B FSV so they cannot accidentally
create or repair a missing database.

For a 1{,}000-event career bundle, the persisted CF row counts
(verified by independent \code{inspect-counts} on every release-C
seed bundle) are:

\begin{center}\small
\begin{tabular}{lr}
\toprule
\code{dj\_raw\_events} & 1{,}004 \\
\code{dj\_adapter\_runs} & 1{,}004 \\
\code{dj\_normalized\_states} & 1{,}000 \\
\code{dj\_instrument\_readings} & 10{,}000 \\
\code{dj\_latent\_panels} & 1{,}000 \\
\code{dj\_transitions} & 1{,}000 \\
\code{dj\_outcomes} & 1{,}000 \\
\code{dj\_dataset\_shards} & 3 \\
\code{dj\_training\_runs} & 1 \\
\code{dj\_model\_artifacts} & 1 \\
\code{dj\_predictions} & 6 \\
\code{dj\_plan\_traces} & 1 \\
\code{dj\_guard\_decisions} & 5 \\
\code{dj\_verification\_runs} & 1 \\
\code{dj\_audit\_log} & 3{,}019 \\
\bottomrule
\end{tabular}
\end{center}

The four-row delta between \code{dj\_raw\_events} and
\code{dj\_normalized\_states} corresponds to four edge-case events
(missing source manifest, invalid action, wrong skill-vector
dimension, out-of-bounds job zone) that are intentionally rejected
by the adapter.

\subsection{Multi-Instrument Latent Panel}

For each declared event, the runtime materializes a
\code{LatentPanel} whose slots are the outputs of declared
\emph{instruments}. An instrument is a frozen function from a typed
event to a structured feature output (a vector, scalar, one-hot, or
pairwise feature). Career-taxonomy uses ten instruments at A1: the
occupation-cluster one-hot, the job-zone scalar, the skill vector,
the action-kind one-hot, the action-effort scalar, the outcome
cluster one-hot, the next skill vector, plus three pairwise
features. The A2 panel adds a transition-distance instrument
(11 slots, input dim 31, target dim 31, latent dim 24); A4 adds
occupation-move-pair pairwise features (12 slots, input dim 51,
target dim 51, latent dim 24, 16{,}688 parameters).

The full ablation panel matrix is in Table~\ref{tab:ablation-panels}.

\begin{table*}[t]
\caption{Ablation panel matrix for \texttt{career\_taxonomy.v1}.
  All numbers read from
  \texttt{paper\_tables/table\_ablation\_multiseed.csv}. ``Action''
  = action conditioning; ``Pairwise'' = pairwise interaction features
  in the panel; ``VICReg'' = variance/covariance regularization
  weights non-zero.}
\label{tab:ablation-panels}
\centering\small
\begin{tabular}{@{}lccccrrrr@{}}
\toprule
\textbf{Ablation} & \textbf{Panel} & \textbf{Action} & \textbf{Pairwise} & \textbf{VICReg} & \textbf{Slots} & \textbf{Input dim} & \textbf{Latent dim} & \textbf{Params} \\
\midrule
A0 & taxonomy only            & yes & no  & yes & 8  & 22 & 24 & 14{,}832 \\
A1 & + skill vector           & yes & no  & yes & 10 & 30 & 24 & 15{,}344 \\
A2 & + transition distance    & yes & no  & yes & 11 & 31 & 24 & 15{,}408 \\
A3 & A2, action removed       & no  & no  & yes & 11 & 31 & 24 & 15{,}408 \\
A4 & A2 + pairwise            & yes & yes & yes & 12 & 51 & 24 & 16{,}688 \\
A5 & A2, no VICReg            & yes & no  & no  & 11 & 31 & 24 & 15{,}408 \\
\bottomrule
\end{tabular}
\end{table*}

\subsection{Predictor and Objective}

The predictor $P_\psi$ is an MLP over the input-panel
representation producing the predicted target panel,
$\hat{s}_{t+1} = P_\psi(s_t, a_t)$. Training minimizes a composite
loss
\begin{align*}
L_{\text{total}} =\ &
  \alpha L_{\text{panel}} + \beta L_{\text{pairwise}}
+ \gamma L_{\text{action}} \\
& + \delta L_{\text{constellation}} + \varepsilon L_{\text{anti-pole}}\\
& + \zeta L_{\text{uncertainty}} + \eta L_{\text{invariant}}
\end{align*}
implemented inside the\linebreak \code{context-graph-dynamicjepa} crate's
\code{train\_dataset} entrypoint, with weights persisted as part of
the model-artifact config so a
reviewer can read them back from \code{dj\_model\_artifacts}. For
the present paper-small results the dominant terms are
$L_{\text{panel}}$ (per-slot MSE/cosine to the frozen target panel)
and $L_{\text{action}}$ (action-conditioning regularizer);
$L_{\text{pairwise}}$ is active only in A4.

\subsection{Planning, Guard, and Surprise}

\code{dynamicjepa plan} enumerates the five candidate actions
(\code{skill\_gain\_adjacent}, \code{skill\_gain\_high\_demand},
\code{role\_move\_near}, \code{experience\_gain},
\code{noop}), runs $P_\psi$ on each, evaluates a constellation /
anti-pole guard, and writes one \code{dj\_plan\_traces} row plus
five \code{dj\_predictions} and five \code{dj\_guard\_decisions}
rows. The selected action is the one whose guarded utility
maximizes the domain-pack-declared scalar.
\code{dynamicjepa record-surprise} compares an observed outcome to
the corresponding prediction; if the surprise exceeds the threshold
a \code{dj\_surprise\_events} row is persisted.

\subsection{Subcommand Surface}
\label{sec:method:cli}

The DynamicJEPA subcommand surface is exactly 29 commands, grouped
in Table~\ref{tab:cli}. The headline command for paper experiments
is \code{research-smoke}, which composes domain registration,
ingest, panel materialization, trajectory and dataset compilation,
training, prediction, planning, ablation runs, edge-case
verification, artifact hashing, paper-table generation, and FSV
evidence emission into one immutable bundle.

\begin{table}[t]
\caption{The 29 DynamicJEPA CLI subcommands.}
\label{tab:cli}
\centering\footnotesize
\begin{tabular}{p{0.4\linewidth}p{0.5\linewidth}}
\toprule
Group & Subcommands \\
\midrule
Domain pack & \code{register-domain}, \code{list-domains},
                \code{get-domain} \\
Inspection & \code{inspect-counts}, \code{inspect-cf} \\
Ingest \& adapt & \code{ingest-event}, \code{run-adapter} \\
Panel & \code{materialize-panel}, \code{get-panel},
                \code{list-instrument-readings},
                \code{bind}, \code{list-bindings} \\
Trajectory \& dataset & \code{compile-trajectories},
                \code{get-trajectory}, \code{list-trajectories},
                \code{compile-dataset}, \code{get-dataset-shard},
                \code{inspect-dataset} \\
Train, predict, plan & \code{train}, \code{get-training-run},
                \code{get-artifact}, \code{predict}, \code{plan},
                \code{record-surprise}, \code{get-prediction},
                \code{get-plan-trace}, \code{get-surprise} \\
Verify per domain & \code{verify-counter-world},
                \code{verify-gridworld},
                \code{verify-career-taxonomy} \\
Research bundle & \code{research-smoke}, \code{hp-sweep},
                \code{multiseed-aggregate} \\
Reference & \code{capture-reference},
                \code{check-reference},
                \code{check-release-}\linebreak[1]\code{aggregate-reference} \\
\bottomrule
\end{tabular}
\end{table}

\subsection{Hardware}
\label{sec:method:hw}

All experiments run on a single host: an NVIDIA GeForce RTX 5090
(compute capability 12.0, driver 596.21, 32{,}607~MiB VRAM, CUDA
13.2 V13.2.78), an AMD Ryzen 9 9950X3D 16-core CPU with 30 logical
threads and 110~GiB of RAM, on Linux 6.6.87.2 (WSL2). The Rust
toolchain is rustc 1.95.0 / cargo 1.95.0. The file
\code{docs2/systemspecs.md}, hashed into every research-smoke
bundle as
\sha{acdc549a5913c59f46376414c08373a8363d646fd3c36d1fc6289a193c0e1205},
records this environment so reviewers can verify the host
description against the bundle.

% ============================================================================
\section{The \code{career\_taxonomy.v1} Public-Taxonomy Benchmark}
\label{sec:bench}

\code{career\_taxonomy.v1} is a public benchmark for action-
conditioned latent transition prediction over occupational data.
The benchmark is deterministically generated from the O*NET 30.2
text database under CC BY 4.0.

\subsection{Source}

The single source is the official O*NET 30.2 text database from
the O*NET Resource Center, fetched once and verified by SHA-256
\sha{b5479271931796b838f7173dc0f673a9ec961b7833ac87168fd11e92e7453741}.
The runtime preflights this archive on every research-smoke run
and fails closed if any of these invariants fail:
\begin{itemize}
\item archive entry count = 42, file count = 41;
\item root prefix = \code{db\_30\_2\_text/};
\item release marker text "O*NET 30.2 Database" present;
\item license marker "Creative Commons Attribution 4.0
      International License" present;
\item per-file row counts: \code{Occupation Data.txt} = 1{,}016,
      \code{Job Zones.txt} = 923, \code{Skills.txt} = 62{,}580,
      \code{Related Occupations.txt} = 18{,}460;
\item per-file required header columns present;
\item no archive member path may escape the archive root.
\end{itemize}
Persisted preflight evidence lives at
\code{onet\_source\_preflight.json} inside every bundle.

\subsection{Walk Generator}

\code{research-smoke --onet-text-zip ...} reads the four required
files, selects the core computer-occupation related-occupation
graph starting at O*NET-SOC \code{15-1252.00} Software Developers,
and walks until the requested number of career events is reached.
Of the 1{,}016 source occupations, 218 are selected; of the
related-occupation edges, 330 are selected. For 1{,}000-event
paper-small runs the generator emits 1{,}000 events; the same
seed-set produces five different deterministic walks under the
same selected edge set. Each generated event's \code{source\_hash}
is bound to the source-ZIP SHA-256, and each transition is
recorded with a per-edge hash in
\code{career\_taxonomy\_onet\_walk.csv} (1{,}001 lines: header plus
1{,}000 transitions). The first generated rows anchor to Software
Developers \code{15-1252.00} and Database Architects
\code{15-1243.00}.

\subsection{Schema}

State has six required fields: \code{occupation\_cluster}
(categorical: \emph{software, data, systems, management}),
\code{job\_zone} (i64, $1{\le}x{\le}5$), \code{skill\_vector}
(vector[4], $0.0{\le}x_i{\le}1.0$), \code{experience\_years}
(f64, $0.0{\le}x{\le}40.0$), \code{wage\_band} (categorical:
\emph{low, mid, high}), and \code{step\_index} (time, i64).

Action is one categorical field \code{move} taking values
\emph{skill\_gain\_adjacent, skill\_gain\_high\_demand,
role\_move\_near, experience\_gain, noop}.

Outcome has four required fields: \code{next\_occupation\_cluster},
\code{next\_job\_zone}, \code{next\_skill\_vector}, and
\code{mobility\_label} (\emph{same, near\_transition,
far\_transition}).

The full TOML is reproduced in Appendix~\ref{app:domainpack}. The
schema choices are deliberately small (4-D skill vector, 4-cluster
occupation, 5-action discrete space) so the entire ablation matrix
fits in a 1{,}000-event paper-small bundle in well under one hour.

\subsection{Splits, Provenance, Edge Cases}

The 1{,}000 generated transitions are deterministically split into
train = 699, val = 200, test = 100 (read from
\code{db\_counts\_after.json} plus the \code{dj\_dataset\_shards}
decoded readback). Splits are persisted in
\code{dj\_dataset\_shards} with leakage checks; no trajectory or
transition appears in more than one shard.

Eleven edge-case rows are generated per bundle and verified
against expected before/after deltas:
\begin{center}\footnotesize
\renewcommand{\arraystretch}{1.1}
\begin{tabular}{@{}lccccc@{}}
\toprule
\textbf{Edge case} & \textbf{Raw} & \textbf{Adp.} & \textbf{State} & \textbf{Trans.} & \textbf{Audit} \\
\midrule
empty input              & 0 & 0 & 0 & 0 & 1 \\
missing source manifest  & 1 & 1 & 0 & 0 & 2 \\
invalid action           & 1 & 1 & 0 & 0 & 2 \\
wrong skill-vector dim   & 1 & 1 & 0 & 0 & 2 \\
out-of-bounds job zone   & 1 & 1 & 0 & 0 & 2 \\
\bottomrule
\end{tabular}
\end{center}
plus six CLI-level edge cases (bad ZIP path, missing
\code{--onet-text-zip} for paper purpose, custom event count
without ZIP, unsafe ZIP member path, incomplete ZIP file count,
too-small \code{--career-event-count}).

\subsection{Reusability}

The benchmark is the smallest reproducible occupational substrate
that requires only the public O*NET ZIP and a recorded SHA-256.
Future work can register an additional domain pack (App.~A) without
changing the runtime, the FSV harness, or the reproducibility
chain.

% ============================================================================
\section{Experiments}
\label{sec:exp}

This section enumerates the experimental matrix. Numerical results
are in \S\ref{sec:results} and read directly from CSVs in the
verified release bundle. Every paper-cited number is derived from
five real O*NET-backed seeds $\{42, 43, 44, 45, 46\}$, each run as
\code{paper\_small\_1k} (1{,}000 career events) and aggregated
through \code{dynamicjepa multiseed-aggregate} with
\code{--min-passed 5} and BCa bootstrap (10{,}000 iterations,
seed \code{20260430}, confidence 0.95).

\subsection{Experiment A --- Runtime Universality}

Does the same runtime, storage, trainer, and FSV harness handle
three unrelated domains by domain pack alone, with no domain-
specific code branches? Domains:
\code{counter\_world.v1} (480 params),
\code{gridworld\_5x5.v1} (16{,}960 params),
\code{career\_taxonomy.v1} (15{,}344 params at A1). Pass requires
that all three domains produce a \code{dj\_verification\_runs} row
with status \code{passed}; that the same five artifact files are
persisted in \code{dj\_model\_artifacts} for each; that all
artifact registry hashes equal recomputed file hashes; and that
the same \code{research-smoke} command produces all three.

\subsection{Experiment B --- Career Next-State Prediction}

Does the trained predictor beat random-init and shuffled-target
baselines on \code{career\_taxonomy.v1}? Headline metrics in
Table~\ref{tab:pred-metrics}.

\subsection{Experiment C --- Action-Conditioned Planning}

Do different candidate actions produce different predicted next-
state panels, and do guard decisions vary with predicted utility?
\code{research-smoke} runs \code{dynamicjepa plan} on a chosen
state with the five candidate actions; five \code{dj\_predictions}
and five \code{dj\_guard\_decisions} rows are persisted; one
\code{dj\_plan\_traces} row records the full search.

\subsection{Experiment D --- Ablation Matrix}

Six ablations A0--A5 (Table~\ref{tab:ablation-panels}). A3 is the
negative control for action conditioning: it persists
\code{model.predictor.ignore\_action=true} in the config; training
and inference \emph{zero} the action vector rather than inventing
a fallback path. The ablation effective seed is the source-fixture
control seed unless \code{--career-ablation-training-seed} is
explicitly supplied, preventing main-seed variation from polluting
ablation comparisons.

\subsection{Experiment E --- Auditability and Failure Behavior}

Does the system fail closed and leave usable diagnostics for every
named edge case? Section~\ref{sec:fsv} catalogs the data-side and
infrastructure-side edges (eleven per career bundle, six aggregate-
level, seven release-level). All emitted typed error codes; none
mutated the output filesystem unless a real bundle was being
created.

\subsection{Experiment F --- Single-GPU Feasibility}

Smoke (120 events) wall time is 169{,}132~ms ($\sim$169~s) for a
full run with 78 commands, 9 verification domains, 11 edge cases,
and 6 ablations. A complete five-seed release-C gate (paper-small
seeds + aggregate + container build + container verify + manual
edges) fits in roughly two hours.

% ============================================================================
\section{Full-State Verification}
\label{sec:fsv}

Full-state verification (FSV) is the central methodological
contribution: every claim in \S\ref{sec:results} is paired with a
physical readback from a \emph{different} code path than the one
that wrote it, and every edge case is verified to fail closed
without mutating the filesystem.

\subsection{The FSV Rule}

For every operation:
\begin{enumerate}[leftmargin=1.2em]
\item Define the source of truth (RocksDB CF rows, file paths,
      table rows).
\item Print before counts and state.
\item Execute the trigger command.
\item Re-read the source of truth through a separate code path
      (\code{inspect-counts}, \code{inspect-cf}, \code{get-*},
      file SHA-256 recomputation).
\item Decode and verify expected values.
\item Run at least three edge cases.
\item Print after counts and state.
\item Write a dated evidence report.
\end{enumerate}

The rule is enforced at three levels: the integration test
\code{dynamicjepa\_research\_smoke\_fsv.rs} (real CLI binary, real
RocksDB, real O*NET ZIP, real CUDA training); the
\code{local\_ci\_gate.sh} script for tier A/B/C runs; and manual
edge-FSV log files under
\code{tmp/mejepa\_local\_ci\_runs/<run\_id>\_manual\_fsv/} with
before / after state printed to disk.

\subsection{The Eight Harnesses}

FSV-1 through FSV-8 cover the runtime's stages: domain pack
registration; instrument registration; event $\to$ panel; binding
and trajectory; predictor dataset; predictor training; planning
and skill execution; guard and surprise. Each has a corresponding
test in \code{crates/context-graph-cli/}\-\code{tests/dynamicjepa\_phase}\linebreak[1]\code{\{3,4,5,6,7,8\}\_fsv.rs}.
The full integration suite passes at \textbf{10 passed, 0 failed,
190.92~s} on the final release-C rerun.

\subsection{Edge-Case Catalog}

Eleven data-shape and CLI-shape edge-case rows per career bundle
(\S\ref{sec:bench}); six aggregate-level edges (existing output
root, count mismatch, duplicate seed, missing required seed, zero
matched bundles, empty seed-failures table); seven release-level
edges (invalid reference kind, missing aggregate path, numeric
drift, existing release output dir, missing Docker tar,
\code{--require-apptainer} without SIF, invalid-SIF-content). All
were manually triggered and produced typed error codes
(\code{VALIDATION}, \code{SCHEMA},
\code{RELEASE\_REFERENCE\_MISMATCH},
\code{RELEASE\_ARCHIVE\_APPTAINER\_INSPECT\_FAILED}, etc.) without
mutating the output filesystem.

\subsection{Read-Only Inspection Hardening}

A subtle defect was found during smoke-B FSV: \code{inspect-counts}
and \code{inspect-cf} previously used the standard writer/open
path, which would \emph{create} a missing RocksDB directory rather
than fail. A read-only inspection tool that creates the source of
truth makes FSV unsound. Fix: the
\code{TeleologicalStoreConfig}'s
\code{create\_missing\_column\_families} flag was added; the
inspection commands now open existing stores only. A hidden \code{--allow-missing} flag exists for explicit
before-state reads on a not-yet-created database, returning
\code{db\_exists=false} with zero counts and not creating the
directory.

\subsection{Hash Reference for the Bundle}

The release-C aggregate bundle is the reference for
\S\ref{sec:results}. Critical SHA-256s in
Table~\ref{tab:aggregate-hashes}. A reviewer can recompute these
locally from the bundle directory and check equality.

\begin{table}[t]
\caption{Aggregate bundle file hashes.}
\label{tab:aggregate-hashes}
\centering\footnotesize
\renewcommand{\arraystretch}{1.15}
\begin{tabularx}{\linewidth}{@{}>{\raggedright\arraybackslash}X >{\raggedright\arraybackslash}X@{}}
\toprule
\textbf{File} & \textbf{SHA-256 (truncated)} \\
\midrule
\sha{aggregate\_manifest.json}             & \sha{734a698c\dots 8670e0} \\
\sha{aggregate\_fsv\_report.json}          & \sha{9b1bc096\dots b6a9b6} \\
\sha{table\_prediction\_metrics\dots.csv}  & \sha{6bd81ee3\dots 02608e7} \\
\sha{table\_ablation\_multiseed.csv}       & \sha{1934260c\dots 936c6a9} \\
\sha{table\_paired\_ablation\_tests.csv}   & \sha{26218e97\dots 88d80} \\
\sha{table\_seed\_failures.csv}            & \sha{8c4e8bfb\dots 5ec52d53} \\
\sha{multiseed\_val\_mse\_box.svg}         & \sha{69db6c4d\dots 23ed77c} \\
\sha{multiseed\_paired\_delta\_forest.svg} & \sha{79f8050a\dots b04e4} \\
\bottomrule
\end{tabularx}
\end{table}

% ============================================================================
\section{Bundles as a Training-Data Class}
\label{sec:training-data}

Up to here this paper has been a verification result. This section
makes a separate claim: the bundles produced by the runtime are not
just a verification artifact. They are also a structured training-
data artifact in their own right, and the class is genuinely new.

\begin{takeaway}
A bundle is mostly deterministic-from-public-source bytes plus an
independent verifier readback. Only a small slice is model-
generated, and that slice is paired with the deterministic observed
target in the same record. Counterfactual ablations and failed runs
are first-class records. The result is a class of training data
distinct from web text, pure synthetic LLM output, and rejection-
sampled chain-of-thought, with substantially lower recursion-
replacement collapse risk.
\end{takeaway}

\subsection{The Synthetic-Data Taxonomy}

The current consensus on training-data quality and model-collapse
risk has four operative categories.
\begin{enumerate}[leftmargin=1.2em]
\item \textbf{Real source data} --- text, images, audio, structured
      records captured from human or natural processes. The
      substrate against which collapse is measured.
\item \textbf{Pure model-generated synthetic data} --- outputs of
      a generative model fed back as training data with no external
      validator. Shumailov et al.\
      \cite{shumailov2023curse,shumailov2024nature} showed that
      this pattern, under recursive replacement, irreversibly
      removes distribution tails.
\item \textbf{Verifier-grounded synthetic data} --- model outputs
      filtered by an external validator (math checker, unit test,
      physics envelope, step-level human label). Process
      supervision and the PRM800K corpus
      \cite{lightman2023verify}, code-execution-grounded
      supervision \cite{jung2025codeexec}, SYNTHETIC-1
      \cite{primeintellect2025synthetic1}, and Enigmata
      \cite{seed2025enigmata} sit here.
\item \textbf{Process-trace data anchored to a deterministic system
      over real source} --- execution traces, code-execution-
      grounded supervision, world-model rollouts, and other
      settings where the bytes of the dataset are deterministic
      functions of a real source plus a runtime, with model outputs
      forming only a small slice of the record and paired with
      independent verification.
\end{enumerate}

The collapse-vs-retain distinction matters: Gerstgrasser et al.\
\cite{gerstgrasser2024accumulate} showed that even with the same
generators, accumulating synthetic data alongside real data avoids
the collapse Shumailov demonstrated for replacement. The pattern
that breaks generative models is not "synthetic" per se; it is
recursive replacement of real distributions by model-induced
distributions.

\subsection{Anatomy of a Bundle as Training Signal}

A research-smoke career bundle decomposes into three byte-slices
with very different provenance.

\textbf{Slice 1 --- Deterministic-from-public-source.}
Generated by walking the O*NET 30.2 related-occupation graph from
\code{15-1252.00} under a fixed seed. These records are
deterministic functions of \linebreak (\code{source\_zip\_sha256},
\code{walk\_seed}, \code{domain\_pack\_hash}):
\code{dj\_raw\_events}, \code{dj\_normalized\_states},
\code{dj\_actions}, \code{dj\_outcomes}, \code{dj\_transitions}
(1{,}004 / 1{,}000 / 1{,}005 / 1{,}000 / 1{,}000 rows for a 1K-event
career bundle); \code{dj\_instrument\_readings} (10{,}000 rows);
\code{dj\_latent\_panels} (1{,}000 rows); \code{dj\_trajectories},
\code{dj\_dataset\_shards} (with leakage checks). A reviewer can
reproduce these byte-for-byte from the public ZIP and the published
walk seed. This slice carries \emph{no} model outputs.

\textbf{Slice 2 --- Tiny-model outputs paired with observed
targets.} For each plan invocation: \code{dj\_predictions} (6 rows
per bundle); \code{dj\_plan\_traces} (1 row, candidate set,
predictions per candidate, guard decisions, utility scores,
selected action, model artifact SHA-256);
\code{dj\_guard\_decisions} (5 rows); \code{dj\_surprise\_events}
(0 or 1 row when observed-vs-predicted delta exceeds threshold).
These \emph{are} model outputs from the 15{,}344-parameter
predictor. Crucially, every prediction lives in the same record as
the deterministic observed next-state from Slice 1, plus the
\code{random\_init\_cosine} and \code{shuffled\_target\_cosine}
baselines tabulated in \S\ref{sec:results}.

\textbf{Slice 3 --- Independent verifier readbacks.}
Persisted through a code path structurally separate from the
writer: \code{dj\_verification\_runs} (1 row per domain);
\code{dj\_audit\_log} ($\sim$3{,}019 rows for a 1K career bundle);
\code{paper\_tables/table\_edge\_cases.csv} (11 typed
\code{(scenario, error\_code, field\_path, remediation)} records);
\code{paper\_tables/table\_ablation*.csv} (6 ablation rows + paired
tests); \code{artifact\_hashes.json} (registry-vs-recomputed file
hashes); \code{onet\_source\_preflight.json} (per-bundle source-ZIP
integrity).

By byte count, Slice 1 $\gg$ Slice 3 $>$ Slice 2 in a typical
career bundle. Predictions, plans, and guards account for tens of
rows; the deterministic walk records and audit log account for
thousands.

\subsection{Properties Relative to Existing Classes}

Table~\ref{tab:tdc-comparison} compares ME-JEPA bundles to the
nearest existing classes. The line item \emph{ME-JEPA bundles
$\neq$ rejection-sampled CoT} is the most counter-intuitive.
Rejection sampling \emph{throws away} the negative samples; the
data the next-generation model sees is the validator-passing slice
only, which is precisely how the distribution narrows over
iterations even with a validator. ME-JEPA bundles persist the
failures (the seed-43 / seed-45 ablation failures, the rc1
missing-binary release-candidate failure, the A3 negative-control
ablation that deliberately collapses to validation cosine 0.790).
Failures are not noise; they are the tail of the distribution
preserved by construction.

\begin{table*}[t]
\caption{Properties of ME-JEPA bundles relative to existing
  training-data classes.}
\label{tab:tdc-comparison}
\centering\scriptsize
\renewcommand{\arraystretch}{1.15}
\setlength{\tabcolsep}{4pt}
\begin{tabular}{@{}p{0.20\linewidth}cccccc@{}}
\toprule
\textbf{Property} & \textbf{Web text} & \textbf{Pure synth.\ CoT} & \textbf{SYNTHETIC-1 / PRM800K} & \textbf{Code-exec traces} & \textbf{World-model rollouts} & \textbf{ME-JEPA bundles} \\
\midrule
Anchored to real source              & yes  & no    & partial    & yes (program) & indirect  & \textbf{yes (O*NET ZIP)} \\
Crypto.\ provenance                  & rare & no    & no         & rare          & no        & \textbf{yes, every row} \\
Per-step verifier                    & no   & no    & yes        & yes           & no        & \textbf{yes (FSV)} \\
Counterfactuals first-class          & no   & some. & no         & no            & some.     & \textbf{yes (A0--A5)} \\
Failed runs preserved                & n/a  & no    & no         & no            & no        & \textbf{yes, by policy} \\
Typed schema per row                 & no   & no    & partial    & yes           & yes       & \textbf{yes (TOML)} \\
Recursion-replace risk               & n/a  & high  & med-low    & low           & medium    & \textbf{low} \\
\bottomrule
\end{tabular}
\end{table*}

\subsection{Collapse-Hazard Analysis}

For the dangerous sense of model collapse --- the
Shumailov / Curse-of-Recursion phenomenon in which recursively
generated content displaces real source data and the resulting
model loses distribution tails --- ME-JEPA bundles have a
substantially lower hazard than typical synthetic data. Three
reasons.

\textbf{Reason 1: Most of the bytes are not model-generated.}
Slice 1 is a deterministic function of the public O*NET ZIP plus
a fixed walk seed. Slice 3 is a separate verifier's readback. Only
Slice 2 is model-generated, and Slice 2 is small. Adding the bundle
to a real-data corpus is closer to \emph{adding a structured
reference dataset} than to \emph{recursive synthesis}.

\textbf{Reason 2: Accumulation, not replacement.} The intended use
is to \emph{accumulate} bundle data alongside existing corpora.
Per Gerstgrasser et al.\ \cite{gerstgrasser2024accumulate},
accumulation does not collapse even with the same generators that
cause collapse under replacement.

\textbf{Reason 3: Tails are explicitly preserved.} The defining
symptom of model collapse is loss of distribution tails. Bundles
preserve tails by three structural mechanisms: failed bundles are
preserved alongside passed runs; eleven edge-case rows per bundle
are typed tail-of-distribution training examples by construction;
and the A3 negative-control ablation forces a low-cosine tail row
into every bundle.

The remaining collapse-hazard slice is Slice 2. Mitigation: every
prediction row is paired with the deterministic observed-next-state
from Slice 1 and the random-init / shuffled-target baselines, so
the natural supervised signal is
\code{(state, observed-next-state)}, with predictions and guards
available as auxiliary process-supervision targets rather than
primary labels.

\subsection{What This Data Class Is Useful For}

Ranked by where it would be worth investing first:

\begin{enumerate}[leftmargin=1.2em]
\item \textbf{Agentic / tool-use models that operate the runtime.}
      Plan traces are exactly what an agent needs to learn the
      \emph{predict $\to$ guard $\to$ act $\to$ log $\to$ audit}
      loop on a typed action space.
\item \textbf{Process-supervised reasoning anchored to a verifier.}
      Lightman-style step-by-step supervision
      \cite{lightman2023verify}, but where the verifier is a
      separate code path that physically re-reads the source of
      truth rather than an LLM-judge.
\item \textbf{Audit / FSV-style reasoning.} Teaching a model to
      always emit before / after counts, cite a SHA-256, and verify
      via a structurally separate path. This pattern is rare in
      current training corpora and would be valuable for any system
      that needs to operate on regulated or auditable data.
\item \textbf{Error-typing and remediation.} The 11 edge-case rows
      plus 6 aggregate-level edges plus 7 release-level edges form
      a small, high-quality corpus of \code{(scenario, typed-error,
      field-path, remediation-string)} quadruples. Almost no public
      dataset has this shape.
\item \textbf{Schema and domain-pack composability.} TOML domain
      packs $\to$ typed records $\to$ typed errors form a high-
      quality corpus for training models that read schemas and
      emit valid records.
\item \textbf{World-model imagination data, scoped.} Plan traces
      are MuZero / Dreamer-style imagined-rollout data on a typed
      action space.
\item \textbf{Code-execution-grounded supervision over a domain
      runtime.} In the spirit of \cite{jung2025codeexec}, but
      grounding supervision in runtime operation rather than
      program execution.
\end{enumerate}

\subsection{What This Data Class Is \emph{Not} Useful For}

A bundle teaches the runtime's behavior on a bounded slice of
O*NET, not "careers." Specifically: not for general world
knowledge (218 selected occupations, 4 clusters, 5 actions); not
for continuous-time or fine-grained temporal dynamics (events are
step-indexed integers); not for modalities outside the declared
instruments (no images, audio, or text-token sequences); not for
open-ended generation (everything is typed). A model trained
naively on a bundle as if it were a general-purpose text corpus
will pick up the schema and the runtime's habits, which is
sometimes useful and sometimes not.

\subsection{Naming the Class}

The proposed working name is \emph{audited execution traces with
cryptographic provenance} (or \emph{audited execution traces} for
short). The defining properties:
\begin{quote}\itshape
Most of the bytes are deterministic functions of a public source
plus a runtime; the small model-output slice is paired in the same
record with an independent verifier readback; counterfactuals and
failures are persisted; every record carries a SHA-256 chain back
to the source.
\end{quote}
The closest existing class is code-execution-grounded supervision
\cite{jung2025codeexec}. ME-JEPA bundles generalize that idea from
program execution to typed domain-pack runtime execution and add
three properties code-execution traces do not routinely carry:
cryptographic provenance to a public source ZIP, independent
verifier readback as a separate record class, and typed
counterfactual ablations as first-class records.

\subsection{Recommendations for Downstream Training}

If a future system uses bundles as part of its training corpus, the
following pattern reduces the risk of importing the runtime's
biases while preserving the data class's strengths:

\begin{enumerate}[leftmargin=1.2em]
\item \textbf{Treat Slice 1 and Slice 3 as primary signal.} They
      are the largest, lowest-risk, highest-provenance portion of
      the bundle.
\item \textbf{Treat Slice 2 as process-supervision targets paired
      with the observed next-state}, not as primary labels. The
      supervised label for next-state is the deterministic observed
      next-state from the same record; predictions and guards are
      auxiliary targets used for process-reward modeling
      \cite{lightman2023verify}.
\item \textbf{Keep failed bundles in the corpus.} They are
      precisely the tail-preserving slice that prevents
      distribution narrowing under the Shumailov pattern.
\item \textbf{Mix with real corpora rather than replacing.} Per
      \cite{gerstgrasser2024accumulate}, accumulation avoids the
      collapse pattern that replacement triggers.
\item \textbf{Bound the claim to the runtime's competence.} A model
      trained on bundle data has learned the runtime's behavior on
      the benchmark's declared schemas. It has not learned the
      world the schemas approximate.
\end{enumerate}

The result is a downstream model that has learned to operate the
audit-aware runtime: predict, plan, guard, log, verify, fail
closed, and emit typed errors with field paths --- a skill set
orthogonal to general world knowledge and currently under-
represented in public training corpora.

% ============================================================================
\section{Results}
\label{sec:results}

All numbers in this section read directly from CSVs in the
multi-seed aggregate bundle's \code{paper\_tables/} directory
inside the release-C run
\sha{20260501T\_release\_c\_gate\_v2}.
Hand-typed numbers are intentionally absent.

\subsection{Headline Multi-Seed Prediction Metrics}

Five seeds $\{42, 43, 44, 45, 46\}$ of \code{paper\_small\_1k}.
BCa bootstrap: 10{,}000 iterations, seed \code{20260430},
confidence 0.95. Career-taxonomy headline:
\begin{itemize}
\item Validation cosine $= 0.905 \pm 0.044$, BCa 95\% CI
      \([0.861, 0.931]\).
\item Test cosine $= 0.868 \pm 0.068$, BCa 95\% CI \([0.806, 0.912]\).
\item Validation action-contrast accuracy
      $= 0.960 \pm 0.043$, BCa 95\% CI \([0.916, 0.982]\).
\item Random-init cosine $= 0.107 \pm 0.255$, BCa 95\% CI
      \([-0.063, 0.370]\).
\item Shuffled-target cosine $= 0.770 \pm 0.107$, BCa 95\% CI
      \([0.671, 0.838]\).
\end{itemize}

The trained predictor's BCa lower bound for validation cosine
$(0.861)$ exceeds the shuffled-target's BCa upper bound
$(0.838)$ by $0.023$ --- the BCa CIs do not overlap. Random-init's
BCa upper $(0.370)$ is far below the trained model's BCa lower
$(0.861)$; this is a clean positive-versus-random comparison.

For the two sanity domains: \code{counter\_world} reaches
validation / test cosine $0.99963$ / $0.99916$ at 480 trainable
parameters; \code{gridworld\_5x5} reaches $0.99473$ / $0.99472$ at
16{,}960 parameters with very tight BCa intervals
(stddev $< 0.0003$ on cosine).

Table~\ref{tab:pred-metrics} reproduces the relevant rows of
\code{table\_prediction\_metrics\_multiseed.csv}.

\begin{table*}[t]
\caption{Multi-seed prediction metrics across three domains.
  Source:
  \texttt{paper\_tables/table\_prediction\_metrics\_multiseed.csv}.
  N=5 seeds; BCa bootstrap 10{,}000 iters, seed
  \texttt{20260430}, confidence 0.95.}
\label{tab:pred-metrics}
\centering\footnotesize
\begin{tabular}{llrrrrrr}
\toprule
Domain & Metric & Mean & Stddev & BCa Low & BCa High & Min & Max \\
\midrule
career\_taxonomy & val\_latent\_mse        & 0.15221 & 0.01027 & 0.14463 & 0.16088 & 0.14072 & 0.16601 \\
career\_taxonomy & val\_latent\_cosine     & 0.90487 & 0.04364 & 0.86150 & 0.93113 & 0.83555 & 0.94767 \\
career\_taxonomy & test\_latent\_mse       & 0.26847 & 0.07205 & 0.22411 & 0.34576 & 0.20173 & 0.38497 \\
career\_taxonomy & test\_latent\_cosine    & 0.86764 & 0.06766 & 0.80630 & 0.91214 & 0.77014 & 0.93346 \\
career\_taxonomy & val\_action\_contrast   & 0.96000 & 0.04345 & 0.91600 & 0.98200 & 0.89500 & 1.00000 \\
career\_taxonomy & random\_init\_cosine    & 0.10703 & 0.25516 & -0.06326 & 0.37000 & -0.15406 & 0.51009 \\
career\_taxonomy & shuffled\_target\_cosine& 0.77046 & 0.10672 & 0.67072 & 0.83775 & 0.61607 & 0.88490 \\
\midrule
counter\_world   & val\_latent\_cosine     & 0.99963 & 0.00000 & 0.99963 & 0.99963 & 0.99963 & 0.99963 \\
counter\_world   & test\_latent\_cosine    & 0.99916 & 0.00000 & 0.99916 & 0.99916 & 0.99916 & 0.99916 \\
counter\_world   & shuffled\_target\_cosine& 0.95389 & 0.00000 & 0.95389 & 0.95389 & 0.95389 & 0.95389 \\
\midrule
gridworld\_5x5   & val\_latent\_cosine     & 0.99473 & 0.00025 & 0.99457 & 0.99497 & 0.99451 & 0.99512 \\
gridworld\_5x5   & test\_latent\_cosine    & 0.99472 & 0.00021 & 0.99459 & 0.99495 & 0.99455 & 0.99507 \\
gridworld\_5x5   & shuffled\_target\_cosine& 0.77882 & 0.01142 & 0.77001 & 0.78734 & 0.76625 & 0.79514 \\
\bottomrule
\end{tabular}
\end{table*}

\subsection{Multi-Seed Ablation Table}

Table~\ref{tab:ablation-results} shows career-taxonomy A0--A5
across the five main career seeds. The ablation training seed is
fixed at the source-fixture control seed for A0, A2, A3, A4, A5
(only A1 varies the main career training seed); cross-seed std for
those rows is therefore zero by design, and the across-ablations
comparison is the informative one.

\begin{table*}[t]
\caption{Multi-seed ablation table (career\_taxonomy). Source:
  \texttt{paper\_tables/table\_ablation\_multiseed.csv}.
  ``Val cos.'' = validation cosine; ``Test cos.'' = test cosine;
  ``Val ACA'' = validation action-contrast accuracy.}
\label{tab:ablation-results}
\centering\small
\begin{tabular}{lrrrrr}
\toprule
Ablation & Params & Val cos.\ (mean$\pm$std) & Test cos.\ (mean$\pm$std) & Val ACA \\
\midrule
A0 (taxonomy only)             & 14{,}832 & $0.97557 \pm 0$        & $0.96301 \pm 0$        & $1.000 \pm 0$ \\
A1 (+ skill vector)            & 15{,}344 & $0.90487 \pm 0.04364$  & $0.86764 \pm 0.06766$  & $0.960 \pm 0.04345$ \\
A2 (+ transition dist.)        & 15{,}408 & $0.95071 \pm 0$        & $0.89696 \pm 0$        & $1.000 \pm 0$ \\
A3 (A2, action removed)        & 15{,}408 & $0.79008 \pm 0$        & $0.65359 \pm 0$        & $0.985 \pm 0$ \\
A4 (A2 + pairwise)             & 16{,}688 & $0.94159 \pm 0$        & $0.92688 \pm 0$        & $1.000 \pm 0$ \\
A5 (A2, no VICReg)             & 15{,}408 & $0.91483 \pm 0$        & $0.81635 \pm 0$        & $1.000 \pm 0$ \\
\bottomrule
\end{tabular}
\end{table*}

\textbf{A3 (action removed) is the predicted negative control.}
Validation cosine collapses from A2's $0.951$ to A3's $0.790$;
test cosine collapses from $0.897$ to $0.654$. A3 still passes its
gate because the floor target was deliberately set to $0.50$ for
the no-action ablation, not because the cosine is high.

\textbf{A0 (taxonomy only) achieves the highest validation cosine.}
This is unintuitive at first glance but consistent with the
diagnostic that the bare taxonomy panel is the easiest-to-fit
target manifold (small, low-dim, mostly categorical). Adding the
skill vector (A1) reduces validation cosine because the 4-D
continuous skill vector is harder to predict than one-hot
occupation cluster. A1 is the headline panel because the test-set
generalization gap (test cosine $= 0.868$) is the most relevant
downstream signal.

\textbf{A4 (pairwise features added) improves test cosine} from
A2's $0.897$ to $0.927$ --- the largest test-set gain in the
matrix --- at the cost of $1{,}280$ additional parameters
($16{,}688$ vs $15{,}408$).

\textbf{A5 (VICReg removed) does not collapse} validation cosine
($0.915$ vs A2's $0.951$) but loses $0.080$ on test cosine
($0.816$ vs A2's $0.897$). The pattern is the predicted one: no
variance / covariance regularization $\Rightarrow$ predictor
degeneracy in the generalization gap.

\subsection{Paired Sign-Flip Permutation Tests}

Source:
\code{paper\_tables/table\_paired\_ablation\_tests.csv}.
Baseline = A0; $n=5$; full sign-flip enumeration over $2^5=32$
sign assignments. The smallest possible two-sided $p$ at $n=5$ is
$2/32 = 0.0625$, which sits above the conventional
$\alpha=0.05$ threshold.
Per the protocol, deltas are reported with BCa CI and explicit
sign-flip $p$ rather than as "significantly different".

The economically meaningful comparison --- A3 vs A0 on validation
cosine --- has mean delta $-0.18549$, BCa CI identically
\([-0.18549, -0.18549]\), sign-flip $p = 0.0625$. The largest delta
in the matrix and the one most directly tied to the design
hypothesis (action conditioning matters).

\subsection{Single-GPU Feasibility}

Smoke (120 events): wall time $169$~s, $78$ commands, $9$
verification domains, $11$ edge cases, $6$ ablations. Paper-small
(1{,}000 events): per-seed runs bounded by training; the
integration FSV suite passes in $190.92$~s. Release-C gate
(5 seeds + aggregate + container build + container verify + manual
edges): $\sim$2 hours wall time. Peak VRAM is far below the
RTX~5090's $32$~GiB --- the trainer's working set never exceeds
$\sim$8~GiB during career-taxonomy training. Headline numbers are
time- and quality-bound, not memory-bound.

\subsection{Failed Runs as Evidence}

The project preserves rather than deletes failed bundles, so each
fixed defect remains traceable. Examples:
\code{20260430T\_paper\_small\_1k\_seed45\_v2} (original seed-45
A3 ablation gate failure, before
\code{--career-ablation-training-seed} separation);
\code{20260430T\_paper\_small\_1k\_multiseed\_42\_46\_v3} (first
all-pass aggregate attempt that failed because empty
\code{table\_seed\_failures.csv} was rejected by a generic CSV
reader);
\code{20260501T\_local\_smoke\_b\_free\_gate\_v3\_hp\_sweep\_batch\_size}
(120-event HP sweep proxy that failed at
\code{batch\_size=32} A2);
\code{mejepa\_5090\_artifact\_v1.0.0\_rc1} (release-candidate v1
missing a packaged verifier binary). All preserved intentionally
as part of the artifact's evidence chain.

% ============================================================================
\section{Reproducibility Artifact}
\label{sec:repro}

The reproducibility artifact is the canonical mechanism for a third
party to regenerate \S\ref{sec:results}'s numbers from a clean
machine. As of 2026-05-01, the artifact is at the local Docker +
Singularity (SIF) release-candidate stage.

\subsection{The Reviewer Path}

\begin{tcolorbox}[colback=gray!7,colframe=gray!50,boxrule=0.4pt,
  left=4pt,right=4pt,top=3pt,bottom=3pt]
\begin{Verbatim}[fontsize=\scriptsize,breaklines=true,breakanywhere=true]
tar -xf mejepa_5090_artifact_v1.0.0_rc5_sif_verified.tar.gz
cd mejepa_5090_artifact_v1.0.0_rc5_sif_verified
singularity exec --nv --cleanenv --no-home \
  -B "$PWD/bundles/.../aggregate:/work/agg:ro" \
  container/contextgraph-mejepa.sif \
  /workspace/reproduce.sh verify-aggregate \
    --aggregate /work/agg \
    --reference /workspace/reference/release.json \
    --json
\end{Verbatim}
\end{tcolorbox}

The Docker invocation is identical except for
\code{docker load -i container/contextgraph-mejepa.docker.tar}
followed by \code{docker run --gpus all ... verify-aggregate}.
Either invocation should print:
\code{\{ status: ok, diff\_count: 0, manifest: passed,
fsv: passed, passed\_count: 5,
rows: \{ ablation\_multiseed\_rows: 78,
paired\_ablation\_test\_rows: 35,
prediction\_multiseed\_rows: 33,
seed\_failure\_rows: 0 \} \}}.

\subsection{Container Strategy}

Two images, same content, two formats. The Docker image is built
on \code{nvidia/cuda:13.2.0-devel-ubuntu24.04} pinned by digest
\sha{sha256:d266e59b88c295bc5fa0e4cef9064eaff84939381b09e2c3d76a5532a303e42d}.
Determinism is checked with two consecutive
\code{docker buildx build} exports using \code{SOURCE\_DATE\_EPOCH}
and \code{type=docker,rewrite-timestamp=true}; both exports must
hash to the same SHA-256. Observed:
\sha{099f0535e9f4f5dc8e07d0afb550ef1e3197657b52ddc7561ef11304f89dd80e}
for both \code{image\_1.tar} and \code{image\_2.tar}.

The Apptainer SIF is built from the Docker tar via the
\code{docker-archive} bootstrap agent
\cite{apptainer2025buildsif} and run with
\code{--cleanenv --no-home --nv}. Resulting SIF: $3.62$~GiB,
SHA-256
\sha{30b531c77966c46524383adcceea9de7983ee251935a68740c2f3009695f2825}.
The packager refuses to include a supplied SIF unless
\code{singularity inspect}, \code{singularity sif list}, and
\code{singularity exec ... verify-aggregate} all pass.

\subsection{Build Determinism}

Every Rust dependency is pinned by \code{Cargo.lock}; the CUDA
base image is pinned by digest, not tag; the binary is built on
the verified host and copied in (no
\code{apt update \&\& apt install} inside the Dockerfile); two
consecutive exports must be byte-identical. There is no claim of
byte-equality across hosts or sm versions: per
\cite{pytorch2026reproducibility}, GPU training can produce small
numeric differences across releases, platforms, or even reruns on
the same host. The release reference therefore uses a
structural-plus-numeric-tolerance schema; cross-run model-weight
bytes are diagnostic evidence, not the pass/fail criterion.

\subsection{Artifact Layout}

The release-candidate archive layout (top-level, abridged):
\code{README.md}, \code{CITATION.cff}, \code{LICENSE} (PolyForm
Noncommercial 1.0.0), \code{THIRD\_PARTY\_LICENSES/ONET.md},
\code{bin/context-graph}
(\sha{41560d160acdaade55f16153678db373c33127ce7a12b4eb5a418d30e86ec0eb}),
\code{container/} (Docker tar + SIF + manifest),
\code{bundles/} (per-seed runs and the aggregate), \code{paper\_tables\_reference/} (the exact CSVs cited in
\S\ref{sec:results}), \code{reference\_hashes.json},
\code{reference\_release\_v1.0.0.json}, \code{reproduce.sh},
\code{manifests/} (LDD readback, SIF inspect, SIF verify),
\code{release\_manifest.json}, \code{archive\_file\_hashes.json},
\code{build\_report.json}, \code{CHANGELOG.md}.

Top-level archive hash:
\sha{8cccf4980893fe949db9b50b8111e592e72ab8419438c4cd08e10c8974943ecd}
for \code{mejepa\_5090\_artifact\_v1.0.0\_rc5\_sif\_verified.tar.gz}.

\subsection{Reviewer Reproducibility Mapping}

Standard reproducibility-checklist questions are pre-filled from
bundle paths: code release (\code{bin/context-graph} +
\code{container/...sif} + \code{container/...docker.tar});
hardware (\code{bundles/<seed>/hardware.json}); software
(\code{environment.json}); seeds (\code{run\_manifest.json});
hyperparameters (\code{training\_metrics.json}); datasets +
licenses (\code{data\_manifest.json}); compute budget
(\code{run\_manifest.json}); limitations (\S\ref{sec:limits}).

% ============================================================================
\section{Discussion}
\label{sec:disc}

\subsection{What Dynamic / ME-JEPA Demonstrates}

\begin{itemize}
\item One Rust binary provides the entire surface; new domains
      ship as TOML.
\item Twenty-five RocksDB column families plus an append-only
      audit log capture every persistent state; readbacks are a
      separate code path.
\item Action-conditioned latent prediction is verifiable end-to-
      end on a single RTX 5090 in under two hours of wall time.
\item Multi-seed BCa CIs and paired sign-flip permutation tests
      are reported uniformly; every paper-cited number reads from
      a CSV in the aggregate bundle, not a hand-typed cell.
\item The runtime fails closed across every named edge case
      without mutating the filesystem (\S\ref{sec:fsv}).
\item Each immutable bundle is itself a structured training-data
      artifact (\S\ref{sec:training-data}). The deterministic
      walk records and verifier readbacks dominate the bytes;
      the small model-output slice is paired in the same record
      with deterministic observed targets and the random-init /
      shuffled-target baselines.
\end{itemize}

\subsection{Failure Modes Observed and Fixed}

Each was found through FSV and persisted as a failed bundle.
\begin{enumerate}[leftmargin=1.2em]
\item Seed-45 / Seed-46 A3 ablation gate failure before
      \code{--career-ablation-training-seed} separation. Fix:
      ablations stay on the source-fixture control seed unless
      explicitly overridden.
\item All-pass aggregate read failure because the empty
      \code{table\_seed\_failures.csv} was rejected by a generic
      CSV reader. Fix: allow empty rows only for that table.
\item Inspection commands silently creating missing RocksDB
      paths. Fix: read-only opens via
      \code{create\_missing\_column\_families = false}; hidden
      \code{--allow-missing} flag for explicit before-state
      reads.
\item HP sweep \code{batch\_size=32} proxy failure at 120 events.
      Fix: HP sweeps use the verified 1K real-O*NET scale and
      scale epochs across batch-size cells.
\item Ruby host verifier inside CUDA container. Fix: verifier
      moved into the compiled Rust CLI
      (\code{check-release-aggregate-reference}).
\item Container determinism failure on Docker iidfile compare.
      Fix: gate compares full deterministic Docker tar exports.
\item Container \code{libcuda.so.1} load failure without
      \code{--gpus all}. Fix: release-C runs container
      verification under the NVIDIA container runtime with
      \code{--gpus all}.
\item Stale release reference. Fix: reference recaptured from a
      preserved passed five-seed aggregate.
\item Initial release-candidate (rc1) missing a packaged verifier
      binary. Fix: packager requires \code{CONTEXT\_GRAPH\_BIN}
      and copies it into \code{bin/context-graph} with a
      recorded \code{manifests/context\_graph\_ldd.txt} dynamic-
      link readback.
\item SIF passed without runtime verification. Fix: packager
      requires \code{singularity inspect --json}, \code{sif list},
      and \code{exec --nv --cleanenv --no-home ... verify-aggregate}.
\end{enumerate}

% ============================================================================
\section{Limitations}
\label{sec:limits}

\subsection{Modeling}

All trained models are tiny: $480$, $16{,}960$, and $15{,}344$
trainable parameters. They are sufficient to demonstrate the
runtime and the FSV protocol but are not state-of-the-art for any
task. Only $H{=}1$ action-conditioned prediction is trained;
multi-step rollouts are a deferred axis. The runtime is latent-
only; top-k occupation prediction and skill-vector reconstruction
require a decoder that is also a deferred axis. All numbers come
from one RTX 5090 host; per
\cite{pytorch2026reproducibility} exact GPU bit-reproducibility
across releases or platforms is not guaranteed, so cross-run
comparisons use tolerance-based comparison.

\subsection{Statistical}

The paired sign-flip permutation test floor at $n=5$ is
$p = 2/32 = 0.0625$, above the conventional $\alpha=0.05$
threshold. Deltas are therefore reported with BCa CI and explicit
sign-flip $p$ rather than as "significantly different". A sample-
size increase to $n=6$ would lower the floor to $p = 0.03125$ and
resolve the threshold.

The cross-seed standard deviation for A0, A2, A3, A4, A5 is zero
because the protocol fixes the ablation training seed at the
source-fixture control seed (only A1 varies the main career
seed). The cross-seed variance reported for those rows is
therefore non-informative about generalization; the
across-ablations comparison is the informative one.

A0--A5 are each compared against A0; A$_i$ vs A$_j$ for
$i,j \neq 0$ are not multiple-tested. The paper-safe consequence
is that no claim is made of "A$_i$ significantly outperforms
A$_j$" for $i,j \neq 0$.

\subsection{Data}

O*NET's coverage reflects U.S. occupational structure; ESCO
cross-walk replication is a deferred axis. The current generator
restricts to core computer occupations rooted at Software
Developers \code{15-1252.00}; this is sufficient for the runtime
contribution but is not coverage-balanced across O*NET's
$1{,}016$ occupations or across O*NET's job-zone distribution.

\subsection{Reproducibility}

The release reference is host-specific: rerunning it on a
different sm\_120 / driver / \code{Cargo.lock} may produce a
slightly different reference hash and is correctly rejected by
the gate. Cross-host portability is future work. The CI gate is
local and self-hosted to keep evidence on disk; a reviewer
without RTX 5090 access cannot run the full release-C gate.

% ============================================================================
\section{Code and Data Availability}
\label{sec:code-data}

The Dynamic / ME-JEPA Rust runtime is currently maintained as a
private codebase. Reviewers and researchers requiring source-level
access for verification or reproduction may request a time-bounded
private grant by writing to the author at
\texttt{chrisroyseai@gmail.com}; reasonable requests from
established research institutions are accommodated.

The canonical reproducible release is the Docker + Apptainer
artifact described in \S\ref{sec:reproducibility} and \S B of the
appendix: a single deterministic Docker tar
(SHA-256 \sha{099f053\dots7d4cd80e}), a single Apptainer SIF
(\sha{30b531c\dots f2825}), the verified release archive
\code{mejepa\_5090\_artifact\_v1.0.0\_rc5\_sif\_verified.tar.gz}
(\sha{8cccf49\dots43ecd}), and the per-seed bundles whose
SHA-chained contents this paper's figures and tables are read
directly from. The release archive is self-contained: a reviewer
with an RTX 5090 host matching the environment fingerprint in
\S\ref{sec:reproducibility} can reproduce every headline number
from the artifact alone, without source-code access.

The career-taxonomy data substrate
(\S\ref{sec:career-benchmark}) derives from O*NET 30.2
\cite{onet2026} under CC BY 4.0; the source ZIP SHA-256 and
per-edge walk hash are recorded in every generated bundle and
are therefore deterministically reproducible from the public
source. A public companion implementation that demonstrates the
same multi-embedder MCP-server architecture in the video
understanding domain is available at
\url{https://github.com/ChrisRoyse/clipcannon} (5 frozen
embedders --- SigLIP-SO400M, Nomic Embed v1.5, Wav2Vec2-large,
WavLM-base-plus-sv, ECAPA-TDNN --- BSL 1.1, converting to
Apache 2.0 on 2030-03-31).

The artifact itself is distributed under PolyForm
Noncommercial 1.0.0; this license choice is intended to
support reproducibility and academic verification while
preserving commercial-licensing optionality, and is not
intended to discourage research use.

% ============================================================================
\section{Acknowledgments}

The author received no external funding for this work and
declares no competing interests.

The runtime builds on prior JEPA work
\cite{assran2023ijepa,bardes2024vjepa,maes2026leworldmodel,chen2026vljepa}
for the latent-prediction premise; on VICReg
\cite{bardes2022vicreg} for the variance / covariance
regularization; on World Models
\cite{ha2018worldmodels}, Dreamer \cite{hafner2023dreamerv3}, and
MuZero \cite{schrittwieser2019muzero} for the planning-with-a-
learned-model framing; and on Perceiver IO
\cite{jaegle2022perceiverio} and Gato \cite{reed2022gato} for the
generality argument. The career-taxonomy substrate uses O*NET 30.2
\cite{onet2026} under CC BY 4.0; ESCO \cite{esco2025} is a planned
cross-walk source. The training-data analysis in
\S\ref{sec:training-data} builds on Shumailov et al.\
\cite{shumailov2023curse, shumailov2024nature} for the model-
collapse phenomenon, Gerstgrasser et al.\
\cite{gerstgrasser2024accumulate} for the accumulation-versus-
replacement distinction, Lightman et al.\
\cite{lightman2023verify} for process supervision, Jung et al.\
\cite{jung2025codeexec} for code-execution-grounded supervision,
and the SYNTHETIC-1 \cite{primeintellect2025synthetic1} and
Enigmata \cite{seed2025enigmata} corpora as reference points. The
reproducibility chain follows Docker reproducible-build guidance
\cite{docker2025reproducible, docker2024bestpractices}, Apptainer
SIF documentation \cite{apptainer2025buildsif}, and MLflow
Tracking \cite{mlflow2024} artifact-lineage patterns. Statistical
reporting follows BCa bootstrap \cite{efron2020automatic} and the
PyTorch reproducibility caveats \cite{pytorch2026reproducibility}.

% ============================================================================
\bibliographystyle{plain}
\bibliography{references}

% ============================================================================
\appendix
\onecolumn
\section*{Appendix}

\subsection*{A. Domain Pack: \texttt{career\_taxonomy.v1}}
\label{app:domainpack}

The full \code{career\_taxonomy.v1} TOML, verbatim from
\code{configs/dynamicjepa/domain\_packs/career\_taxonomy.v1.toml}:
\begin{tcolorbox}[colback=gray!4,colframe=gray!40,boxrule=0.4pt,
  left=3pt,right=3pt,top=2pt,bottom=2pt]
{\footnotesize
\begin{verbatim}
[domain]
id = "career_taxonomy"
version = "1.0.0"
title = "Career Taxonomy - Public Occupational Transition Smoke Benchmark"
schema_version = 1

[time]
field = { name = "step_index", kind = "i64", required = true, min = 0, max = 1000000 }

[[state.fields]]
name = "occupation_cluster"
kind = { type = "categorical", variants = ["software", "data", "systems", "management"] }
required = true

[[state.fields]]
name = "job_zone"
kind = { type = "i64" }
required = true
min = 1
max = 5

[[state.fields]]
name = "skill_vector"
kind = { type = "vector", dim = 4 }
required = true
min = 0.0
max = 1.0

[[state.fields]]
name = "experience_years"
kind = { type = "f64" }
required = true
min = 0.0
max = 40.0

[[state.fields]]
name = "wage_band"
kind = { type = "categorical", variants = ["low", "mid", "high"] }
required = true

[[action.fields]]
name = "move"
kind = { type = "categorical", variants =
  ["skill_gain_adjacent", "skill_gain_high_demand",
   "role_move_near", "experience_gain", "noop"] }
required = true

[[outcome.fields]]
name = "next_occupation_cluster"
kind = { type = "categorical", variants =
  ["software", "data", "systems", "management"] }
required = true

[[outcome.fields]]
name = "next_job_zone"
kind = { type = "i64" }
required = true
min = 1
max = 5

[[outcome.fields]]
name = "next_skill_vector"
kind = { type = "vector", dim = 4 }
required = true
min = 0.0
max = 1.0

[[outcome.fields]]
name = "mobility_label"
kind = { type = "categorical", variants =
  ["same", "near_transition", "far_transition"] }
required = true

[[entity.fields]]
name = "trajectory_id"
kind = { type = "string" }
required = true

[[entity.fields]]
name = "source_dataset"
kind = { type = "string" }
required = true
\end{verbatim}
}
\end{tcolorbox}

\subsection*{B. Hash Reference}

\begin{itemize}\small
\item O*NET 30.2 source ZIP\\
  \sha{b5479271931796b838f7173dc0f673a9ec961b7833ac87168fd11e92e7453741}
\item Aggregate \code{aggregate\_manifest.json}\\
  \sha{734a698c7ed946c292cf9005c4083f7205a2a37b271037870418103eb68670e0}
\item Aggregate \code{aggregate\_fsv\_report.json}\\
  \sha{9b1bc096c850fbd87d44447f53d33da58c35cd4337292ffbde82f63cc6b6a9b6}
\item \code{table\_prediction\_metrics\_multiseed.csv}\\
  \sha{6bd81ee3b53b90c495263378c492a7fed2b123cbcfd5007ee0a14a99702608e7}
\item \code{table\_ablation\_multiseed.csv}\\
  \sha{1934260ce6d0a39b6d04d0c857e097ff3d19c3e407fca6b00760548e5936c6a9}
\item \code{table\_paired\_ablation\_tests.csv}\\
  \sha{26218e9705c2170b25a7251f51234a4fb18b02eed9ab92339189f22a6dd88d80}
\item \code{table\_seed\_failures.csv}\\
  \sha{8c4e8bfb875f646993fb453e8e35203faca6fd99c41373ed8ef1ed455ec52d53}
\item Docker tar (deterministic, two-of-two)\\
  \sha{099f0535e9f4f5dc8e07d0afb550ef1e3197657b52ddc7561ef11304f89dd80e}
\item Apptainer SIF\\
  \sha{30b531c77966c46524383adcceea9de7983ee251935a68740c2f3009695f2825}
\item CUDA base image (pinned by digest)\\
  \sha{sha256:d266e59b88c295bc5fa0e4cef9064eaff84939381b09e2c3d76a5532a303e42d}
\item Release archive
  \code{mejepa\_5090\_artifact\_v1.0.0\_rc5\_sif\_verified.tar.gz}\\
  \sha{8cccf4980893fe949db9b50b8111e592e72ab8419438c4cd08e10c8974943ecd}
\item \code{bin/context-graph}\\
  \sha{41560d160acdaade55f16153678db373c33127ce7a12b4eb5a418d30e86ec0eb}
\item Host environment fingerprint
  (\code{docs2/systemspecs.md})\\
  \sha{acdc549a5913c59f46376414c08373a8363d646fd3c36d1fc6289a193c0e1205}
\end{itemize}

\subsection*{C. Glossary}

\textbf{JEPA.} Joint-Embedding Predictive Architecture. Predicts
target representations rather than reconstructing observations.

\textbf{ME-JEPA.} Multi-Embedder JEPA. Targets are produced by a
panel of frozen, externally-defined per-domain instruments.

\textbf{Dynamic / ME-JEPA.} The runtime in this paper. Adds
domain packs, source-of-truth storage, full-state verification,
and audited planning / surprise.

\textbf{Domain pack.} A TOML manifest declaring state schema,
action schema, instruments, and panel slots; the only thing a new
domain ships.

\textbf{FSV.} Full-State Verification. Every operation has a
separate read path that re-reads the source of truth.

\textbf{CF.} RocksDB column family. The runtime uses 25
DynamicJEPA-prefixed CFs plus the audit log.

\textbf{BCa bootstrap.} Bias-corrected and accelerated bootstrap
confidence interval. 10{,}000 iterations, seed \code{20260430},
confidence 0.95.

\textbf{Sign-flip permutation test.} Paired-sample non-parametric
test with two-sided $p$ obtained by enumerating all $2^n$ sign
assignments ($p = 0.0625$ floor at $n=5$).

\textbf{Audited execution traces.} The training-data class
introduced in \S\ref{sec:training-data}. Most bytes are
deterministic functions of a public source plus a runtime; the
small model-output slice is paired with deterministic observed
targets and external baselines; counterfactuals and failures are
persisted; every record carries a SHA-256 chain back to source.

\textbf{Slice 1 / Slice 2 / Slice 3.} The three byte-slices of a
bundle for training-data purposes (\S\ref{sec:training-data}):
deterministic-from-source, tiny-model-output paired with observed
targets, and independent verifier readbacks, respectively.

\textbf{Model collapse.} The Shumailov-et-al.\ phenomenon in which
recursively-replaced model-generated training data causes
irreversible loss of distribution tails.

\textbf{Process supervision.} Per-step labels for reasoning
chains, as opposed to outcome-only labels.

\textbf{Release candidate (rc5).} The current local Docker + SIF
release-candidate archive.

\end{document}
